% Nouveau EGA V (cote 25): archive cote 25, site page 16 (PDF page 17).
% typescript. Draft compared: none.
% AI-transcribed by Claude Opus 5.5 (claude-opus-5-5, adaptive thinking, effort high), Message Batches API, batch ega25b (msgbatch_01A7n5YuXubUPNK1zsepi8RG), 2026-11-13; not yet checked by eye.
\scanpage{25/16}
\mspage{\hand{13}}

\var{\deleted{\uline{Corollaire} 2.8. Sous les conditions de 2.7. , pour que $Z_i$ soit non
immergé, il f et s que $Z_i$ \var{\inserted[typed]{sur \unsure{$\bar{k}$}}}{sur \unsure{$\bar{k}$}}{insertion}{ega25_25_016-v0002} soit non immergé. S'il en est ainsi, la
multiplicité radicielle de $F$ en $Z_i$ est égale à la multiplicité radici-%
elle de $G$ sur $k(\ \ )$ en $Z_i$ .

La première assertion résulte aussitôt de la dernière assertion
de 2.7. Quant à la première, on peut supposer\var{\xout{, (comme il s'agit d'une}}{}{typewriter strike-out}{ega25_25_016-v0003}
\var{\inserted[typed]{par restriction à un ouvert de X et}}{par restriction à un ouvert de X et}{insertion}{ega25_25_016-v0004}
\var{\xout{\unsure{notion géométrique}}}{}{typewriter strike-out}{ega25_25_016-v0005} par dévissage sur $F$ (comme $Y$ est plat sur $X$
donc $F\ \ G$ est une foncteur exact) que l'on a $F=\underline{O}_{Z_i}$ , et on peut
évidemment remplacer $X$ par $Z_i$ . On est donc ramené au cas où $X$ est intè-%
gre et $F=\underline{O}_X$.}}{}{deletion}{ega25_25_016-v0001} Soit $K/k$ le\var{\xout{s}}{}{typewriter strike-out}{ega25_25_016-v0006} corps des fonctions de $X$, et soit $n=\dim P$;
introduisant des coordonnées \var{\xout{projec}}{}{typewriter strike-out}{ega25_25_016-v0007} affines $T_1,..,T_n$ dans $P$ (en choisiss
ant un hyperplan à l'infini $H_\hand{\infty}$ tel que $f(X)$ ne soit pas contenu dans
celui-ci) et $S_1,...S_n$ dans $\overset{\hand{\vee}}{P}$ , on voit que le corps\var{\xout{\unsure{ex}}}{}{typewriter strike-out}{ega25_25_016-v0008}\var{\inserted[typed]{des}}{ des}{insertion}{ega25_25_016-v0009} fonctions $L$
de $Y_{\hand{\eta}}$ \var{\xout{xxx}}{}{typewriter strike-out}{ega25_25_016-v0010} s'identifie au corps des fractions de l'anneau intègre
$K[S_1,..,S_n]/(\hand{\sum} t_iS_i - 1)$ , où les $t_i \hand{\in} K$ sont les images des $T_i$ à
l'adde de $f:X\hand{\to} P$ . Comme $\dim \hand{\overline{\textcolor{black}{f(X)}}} \hand{>} 0$ , les $t_i$ ne sont pas tous
algébriques sur $k$, à fortiori ils ne sont pas tous nuls, supposons
p. ex. $t_n \hand{\neq} 0$. On constate alors aussitôt qu'on a $L = K(S_1,..,S_{n-1})$ ,
(extension transcendante pure), $S_n\hand{\in} L$ étant donné par l'équation
\[ \hand{\sum} t_iS_i - 1 = 0 \]
en fonction des $S_i$ $(1\hand{\leq} i\hand{\leq} n-1)$ et des $t_i$ $(1\hand{\leq} i\hand{\leq} n)$. D'autre part,
$k' = k(\hand{\eta})$ s'identifie à $k(S_1,..,S_n)$, et l\var{\xout{a}}{}{typewriter strike-out}{ega25_25_016-v0011}'inclusion canonique
$k'\hand{\longrightarrow} L$ s'obtient en envoyant $S_i$ sur $S_i$ , i.e. $k'$ en tant que sous-%
\var{\xout{\unsure{exten}}}{}{typewriter strike-out}{ega25_25_016-v0012}\var{\inserted[typed]{extension}}{extension}{insertion}{ega25_25_016-v0013} de $L$ est la sous-extension engendrée par les $S_i$ $(1\hand{\leq} i\hand{\leq} n)$, ou
ce qui revient évidemment au même, par les $S_i$ $(1\hand{\leq} i\hand{\leq} \var{\xout{n}}{}{typewriter strike-out}{ega25_25_016-v0014}\var{\inserted[typed]{n-1}}{n-1}{insertion}{ega25_25_016-v0015})$ et par
$S_n = a_o + a_1S_1 + \ldots + a_{n-1}S_{n-1}$ , où $a_o = t_n^{-1}$ , $a_i =-t_it_n^{-1}$ pour
$1\hand{\leq} i\hand{\leq} n-1$. \var{\deleted{Comme}}{}{deletion}{ega25_25_016-v0016}\var{\inserted{Notons que}}{Notons que}{insertion}{ega25_25_016-v0017} le corps engendré par les $a_i$ ou par les $t_i$ est évidem
ment le même, \var{\deleted{les $a_i$ \unsure{ne|n} sont également pas tous algébriques sur k.}}{}{deletion}{ega25_25_016-v0018}\var{\inserted{leur degré de transcendance commun
n'est autre que $\dim f(X)$.}}{leur degré de transcendance commun
n'est autre que $\dim f(X)$.}{insertion}{ega25_25_016-v0019}

\var{\margin{16}}{\placedmargin{16}}{marginal note placed}{ega25_25_016-v0020}
